\scanpage{5}
UNIVERSITE\\
DES SCIENCES ET TECHNIQUES\\
DU LANGUEDOC\\
\texttt{------------}\\
Service de la Scolarité\\
MA/GJ\\
Diffusion : Liste A\\
N\textdegree\ : 5890

Montpellier, le 7 Novembre 1980

AVIS DE SOUTENANCE D'UNE THESE DE DOCTORAT DE 3ème CYCLE

ECOLOGIE GENERALE ET APPLIQUEE Opt. ECOLOGIE TERRESTRE

\figure{\begin{tikzpicture}\draw (0,0)--(4,0);\end{tikzpicture}}

\uline{Nom et prénom du candidat} : MICHELLAND Dominique

\uline{Titre de la thèse} : SURVIE EN MILIEU INSULAIRE : QUELLE STRATEGIE ? LE CAS DES MESANGES EN CORSE.

\uline{Date, heure et lieu de soutenance prévus} :\\
17 NOVEMBRE 1980 à 15 Heures\\
Salle de cours n\textdegree\ 9 - USTL

\uline{Nom des Directeurs de thèse} :\\
MM. LEGENDRE et BLONDEL

\uline{Désignation du Laboratoire où la thèse est préparée} :\\
Laboratoire de Systématique et d'Ecologie Méditerranéennes

\uline{Résumé de la thèse} :

En Corse les populations de mésanges présentent un certain nombre de caractères démographiques et biométriques qui les différencient nettement des populations continentales.

A la lumière de la théorie de la biogéographie insulaire, cet ensemble d'observations peut s'interpréter comme le résultat de pressions de sélection de type K qui favorisent la survie des populations dans les milieux isolés et confinés.

Ce travail montre les répercussions, à l'échelle des populations, des importantes modifications, mises en évidence par ailleurs, de la structure et de la dynamique des communautés insulaires.

\uline{Durée approximative des recherches ayant abouti à la thèse} : 2 ans

\uline{Composition du jury proposé} :\\
MM. AMANIEU\\
\hspace*{2em}BLONDEL\\
\hspace*{2em}BOURLIERE\\
\hspace*{2em}THALER\\
\hspace*{2em}ISENMANN

P. le Secrétaire Général,\\
l' Attachée d'Administration Universitaire,

\uline{\unsure{P. Gentil|A. Gentil}}\\
\uline{A. GENTIL}

\figure{\begin{tikzpicture}
\draw (0.3,2.5) arc[start angle=100,end angle=245,radius=1.55];
\draw[gray] (-0.97,1.8) .. controls (-1.2,1.2) and (-0.8,0.3) .. (-0.6,0);
\node[rotate=62] at (-0.95,1.95) {\unsure{ACAD}};
\node[rotate=-41] at (0.05,1.05) {\unsure{SCIENCES}};
\end{tikzpicture}}

\mspage{4}
